\section{Reproducibility}\label{app:repro}
\setcounter{table}{0}\setcounter{figure}{0}\setcounter{equation}{0}

The runs whose reports are committed to the repository are listed below, by the names of
\cref{tab:runs}, with their configurations and records. Each configuration hash is the
SHA-256 of the configuration's canonical JSON form (keys sorted, no whitespace), recomputed
from the committed file; it agrees with the hash that the run recorded at start-up. The
initial three-dimensional run and the pilot run of its amendment are documented in the
repository's list of deviations, from which the values quoted for them are taken.

% generated by scripts/make_cmame_material.py from the records; do not edit
\begingroup\raggedright
\begin{itemize}
\item July run (two dimensions), 16 July 2026: configuration \path{configs/phase1_rec8_v2.json} (SHA-256 \texttt{62b26ad868d424ef}\ldots), record \path{records/phase1/report_rec8_v2.json} (SHA-256 \texttt{c6ea5934ff7fee4d}\ldots).
\item Diagnostic run (two dimensions, exploratory), 30 July 2026: configuration \path{configs/diag_anchor_b16.json} (SHA-256 \texttt{1a85eeabf2646568}\ldots), record \path{records/phase1/report_diag.json} (SHA-256 \texttt{dbd7e13dc28dea4b}\ldots).
\item Joint-embedding run (two dimensions), 31 July 2026: configuration \path{configs/wp2_e2_v2.json} (SHA-256 \texttt{60088135b1fecb27}\ldots), record \path{records/phase1/report_wp2_e2.json} (SHA-256 \texttt{1a868e9ca6b72d58}\ldots).
\item Three-dimensional comparison (the amended run), 28 September 2026: configuration \path{configs/phase2b_v1.json} (SHA-256 \texttt{316f5e6e282db9c1}\ldots), record \path{records/wp8/e2/baseline/report_phase2b.json} (SHA-256 \texttt{320b6db5060ecae9}\ldots).
\item Latent-regularisation run, label-free network, 29 September 2026: configuration \path{configs/e1_2d_base.json} (SHA-256 \texttt{4dfdea42e22c0cc4}\ldots), record \path{records/wp8/e1/e1_2d_base/report.json} (SHA-256 \texttt{0ebd53c3be1dcaf8}\ldots).
\item Latent-regularisation run, regularised network, 29 September 2026: configuration \path{configs/e1_2d_shaped.json} (SHA-256 \texttt{9433e7cb60594be7}\ldots), record \path{records/wp8/e1/e1_2d_shaped/report.json} (SHA-256 \texttt{a644f899502a0180}\ldots).
\item Latent-regularisation run, regulariser on the latent vectors directly, 30 September 2026: configuration \path{configs/e1_2d_raw_s0.json} (SHA-256 \texttt{aa1d44d28563e7c6}\ldots), record \path{records/wp8/e1/e1_2d_raw_s0/report.json} (SHA-256 \texttt{e46664a472c5fbf4}\ldots).
\item Bottleneck run, M = 512, 30 September 2026: configuration \path{configs/e2_m512.json} (SHA-256 \texttt{82118d7d3309465d}\ldots), record \path{records/wp8/e2/e2_m512/report.json} (SHA-256 \texttt{02db3bb6715d7886}\ldots).
\item Bottleneck run, M = 1,024, 30 September 2026: configuration \path{configs/e2_m1024.json} (SHA-256 \texttt{526a5688eb49f333}\ldots), record \path{records/wp8/e2/e2_m1024/report.json} (SHA-256 \texttt{3bcd56f14c74ac2a}\ldots).
\item Two-dimensional comparison, 9 October 2026: configuration \path{configs/cm2d_v1.json} (SHA-256 \texttt{bf1f1143b397331b}\ldots), record \path{records/cmame/cm2d/return/report.json} (SHA-256 \texttt{a72c007e1a612986}\ldots).
\end{itemize}
\par\endgroup


Every number in the text that reports a result of this study, every table except
\cref{tab:runs}, and every figure but the schematic of \cref{fig:schematic} are generated
by \path{scripts/make_cmame_material.py} from these records, from the readings taken after
the runs and their benchmark records, from the timings of \cref{tab:cost}, and from the
list of deviations, with the bound of the
reuse check read from the adjudicator of the two-dimensional comparison, whose hash its
verdict records; the inputs and their SHA-256 hashes are listed in
\path{paper/cmame/generated/sources.json}. The readings taken after the runs are those of
the amplitude (\cref{sec:transfer:amplitude}), of the error spectra of the two-dimensional
comparison (\path{scripts/cm2d_spectra.py}; \cref{sec:results2d:spectrum}) and of the
per-load-case energies, error norms and fields of the three-dimensional networks
(\path{scripts/export_fields.py}; \cref{sec:results3d}); each was run on the training
machine on the stored networks, whose hashes it records, and is committed with the record
of its run. No result of this study is typed by hand. The design constants stated in the
text, and the values quoted from other studies with their citations, are typed; the tests
check the design constants against the configurations and the code. The test file
\path{tests/test_cmame_material.py} regenerates the material and compares it with the
committed files byte for byte; checks the hashes of the two-dimensional records; recomputes
the main numbers of the text independently from the records, and every reading of the
spectra and of the fields from their per-load-case arrays; checks the network sizes against
the code; and checks \cref{prop:stress,prop:amp}, \cref{cor:rank,cor:zero,cor:warm} and
\cref{rem:rank,rem:residual} on assembled two- and three-dimensional instances. The
generator itself checks the statements that the text makes about the records, for example
that a network was more accurate on every metric in each seed, and stops if a record
disagrees. The mean and standard deviation over seeds in the tables are the mean and the
sample standard deviation of the per-seed values, which differs from the population
standard deviation stored in the run reports. The code was written with the assistance of
Claude Fable 5, Claude Fable 5.1 and Claude Opus 5.5 (Anthropic), as declared at the end of
the paper; the tests described above check it. The schematic of \cref{fig:schematic} and the
graphical abstract are drawn by TikZ code written with the assistance of Claude. Every other
figure is drawn by the generator script from the committed records; no image was produced
or altered by a generative AI tool.

The hosting service's record of the repository's activity, which logs the creation of each
branch and each push to it, including the force push described in \ref{app:prereg}, was
exported on 10 October 2026 and is committed in \path{records/provenance/} with a note on
how it was obtained. Beside it are the entries that name the repository in the security
log of the account that created it, exported the same day, with their user agents and
request identifiers removed. On 11 October 2026 the repository was transferred from that
account to the first author's and renamed. Its record of activity, exported again after the
transfer and committed beside the first export, begins with the 103 events of the first,
and Software Heritage, which archived the repository from its former address before the
transfer and from its present address after it, recorded the same snapshot both times, the
one given in the Data availability statement.

The repository's address and the tag of the version described here are given in the Data
availability statement, with the licences.

\subsection{Executable checks}\label{app:repro:checks}

The statements of \cref{sec:theory} are checked numerically, as follows. \Cref{lem:exact}
holds on assembled instances to machine precision. The checks of \cref{lem:cond},
\cref{cor:modes} and \eqref{eq:cheb} were pre-registered for sixteen assembled validation
instances of each corpus (criteria C5 and KP5 in \cref{tab:criteria}); none was violated.
The modewise identity \eqref{eq:modes} held to \numModesPhrase{} on the extreme
eigenvectors; for random perturbations of the solutions, the displacement errors reached at
most \numCondTwo{} and \numCondThree{} of the bound \eqref{eq:cond}, and the
conjugate-gradient ratios at most \numChebTwo{} and \numChebThree{} of the bound
\eqref{eq:cheb}; the condition numbers of these instances are given in
\cref{sec:methods:data}. On the \numAmpRows{} instance-seed evaluations of the readings of
\cref{sec:transfer:amplitude}, the relative energy gap after rescaling, computed from the
energies, agreed with $1-\cos^2\psi$ to within \numAmpIdentity{}. On every validation load
case of the networks measured after the runs, all with 1,024 labelled instances (the four
trained networks of the two-dimensional comparison, and the label-free transformer, the
supervised transformer and the graph network in three dimensions), the energy was positive
exactly when the error's stiffness norm exceeded the solution's, that is, when the
prediction was further from the solution than the zero field in the stiffness norm
(\cref{cor:zero}), the rescaling by $\cstar$ never increased the relative energy gap
(\cref{prop:amp}), and the area- or volume-weighted von Mises error obeyed
\eqref{eq:vmboundrel}, with largest ratios of its left side to its right side of
\numSpecBoundMax{} in two dimensions and \numFieldBoundMax{} in three.

The area- or volume-weighted von Mises errors of these networks differed little from the
element-wise ones that the tables report (\cref{sec:theory:stress}). In two dimensions
their medians over the validation load cases were \numSpecVmAreaFree{},
\numSpecVmAreaKnorm{}, \numSpecVmAreaLab{} and \numSpecVmAreaMgn{} for the label-free, the
stiffness-norm and the supervised transformers and the graph network, against element-wise
medians of \numSpecVmElemFree{}, \numSpecVmElemKnorm{}, \numSpecVmElemLab{} and
\numSpecVmElemMgn{}; in three dimensions they were \numFieldVmVolFree{},
\numFieldVmVolLab{} and \numFieldVmVolMgn{} for the label-free transformer, the supervised
transformer and the graph network, against element-wise medians of \numFieldVmElemFree{},
\numFieldVmElemLab{} and \numFieldVmElemMgn{}.

\subsection{TF32 rounding}\label{app:repro:tf32}

The networks ran with TF32 matrix products (\cref{sec:methods:training}), whose rounding
varies from node to node and therefore weighs heavily in the stiffness norm, so each
network of the two-dimensional comparison was also evaluated with TF32 off, after the run.
The rounding, the difference between the two predictions, had a squared stiffness norm of a
median \numSpecTfFree{} and \numSpecTfKnorm{} of that of the error for the label-free and
the stiffness-norm transformers (\numSpecTfLab{} and \numSpecTfMgn{} for the two networks
trained on the displacement error). In the modes with $\lambda_m\ge10^4R(\Ustar)$ it had a
median \numSpecTfStiffFree{} and \numSpecTfStiffKnorm{} of the error's squared stiffness
norm in those modes, on the \numSpecTfStiffN{} triples whose spectra reach that far
(\numSpecTfStiffLoads{} of the \numSpecLoadCases{} pairs of instance and load case, in each
seed). With TF32 off, the median quotients $\rho$ were \numSpecIeeeRqFree{} (label-free),
\numSpecIeeeRqKnorm{} (stiffness norm), \numSpecIeeeRqLab{} (supervised) and
\numSpecIeeeRqMgn{} (graph network), and the geometric-mean ratios of $\rho$ and of $g$
between the supervised and the stiffness-norm transformers rose to
\numSpecIeeeLabKnormRq{} and \numSpecIeeeLabKnormGap{}. The rounding thus places the errors
of the two networks trained in the stiffness norm higher in the spectrum than they would lie
without it, and the readings as run understate the differences between the networks trained
in the two norms.

\subsection{Field plots}\label{app:repro:plots}

The instances shown in the field plots of \cref{sec:results3d,sec:results2d,sec:transfer}
were chosen by rules fixed in code before any field was drawn: in three dimensions, the
validation instance with the largest relative energy gap of the supervised transformer (the
same instance in all three seeds), the validation instance with the median relative energy
gap of the label-free transformer, and the fine-mesh instance with the median displacement
error of the label-free transformer; in two dimensions, the validation instance with the
median relative energy gap of the supervised transformer; each in seed 0, the median of $n$
values being the lower one, of rank $\lfloor(n-1)/2\rfloor$ counted from zero in ascending
order. The first three-dimensional rule, and the two-dimensional rule, which was written
after the verdict of the two-dimensional comparison, were chosen after the per-instance
records had been examined, but before any field had been plotted. Each plot shows one load
case of its instance, chosen by the same rule applied to the instance's four load cases (the
largest value, or the second smallest); this was fixed when the plots were made, after the
per-load-case values had been examined but before any field had been drawn. The
three-dimensional plots show the three faces of the box seen in an oblique projection, the
free end on the right; the clamped face and the cavities are hidden.
